\documentclass[11pt]{article}

\usepackage[T1]{fontenc}
\usepackage{lmodern}
\usepackage[a4paper,margin=28mm]{geometry}
\usepackage{amsmath,amssymb,amsthm,mathtools}
\usepackage{microtype}
\usepackage{enumitem}
\usepackage{xcolor}
\usepackage[colorlinks=true,linkcolor=blue!45!black,citecolor=blue!45!black,
  urlcolor=blue!55!black]{hyperref}
\setlength{\emergencystretch}{2em}

\newtheorem{theorem}{Theorem}[section]
\newtheorem{proposition}[theorem]{Proposition}
\newtheorem{lemma}[theorem]{Lemma}
\newtheorem{corollary}[theorem]{Corollary}
\newtheorem{definition}[theorem]{Definition}
\newtheorem{remark}[theorem]{Remark}

\newcommand{\C}{\mathbb C}
\newcommand{\R}{\mathbb R}
\newcommand{\N}{\mathbb N}
\newcommand{\D}{\mathbb D}
\newcommand{\cA}{\mathcal A}
\newcommand{\cG}{\mathcal G}
\newcommand{\cH}{\mathcal H}
\newcommand{\cM}{\mathcal M}
\newcommand{\cP}{\mathcal P}
\newcommand{\cS}{\mathcal S}
\newcommand{\cT}{\mathcal T}
\newcommand{\FE}{\mathsf R}
\newcommand{\Span}{\operatorname{span}}
\newcommand{\LogDer}{\operatorname{ld}}
\newcommand{\cl}{\operatorname{cl}}

\title{Cofinite Universality of Reflected Lambert Spectra\\
and a Local-Source Barrier for the Riemann Hypothesis}
\author{Michael Haufschild}
\date{}

\begin{document}

\maketitle

\begin{abstract}
We prove a spectrum-independent completeness law for reflected derivative-Lambert
kernels and use it to identify an exact obstruction to local-source approaches to the
Riemann hypothesis.

Let \(\Lambda=(\lambda_j)\) be an injective positive frequency family.  Assume that for
every integer \(n\geq1\), every sufficiently large real \(R\), and after any finite
deletion,
\[
 \sum_{\lambda_j>R+1}\frac{1}{(\lambda_j-R)^n+1}=\infty.
\]
On every closed disc centred at \(1/2\) and contained in \(0<\Re s<1\), the real span of
\[
 \lambda_j\frac{e^{-\lambda_j s}}{(1-e^{-\lambda_j s})^2}
 -
 \lambda_j\frac{e^{-\lambda_j(1-s)}}{(1-e^{-\lambda_j(1-s)})^2}
\]
is dense in the full space of holomorphic, functional-equation-odd,
Schwarz-real germs.  More strongly, the additive group generated by arbitrarily
small, pole-free multiplicative local-factor increments has the entire germ fibre as
its orbit closure.  Consequently its maximal continuous Hausdorff quotient is one
point.  The result is cofinite and does not depend on arithmetic.

For \(\lambda_p=\log p\), the hypothesis follows from Euler's divergence of
\(\sum_p p^{-1}\), with no prime number theorem.  We then couple the orbit law to the
normalized prime--archimedean logarithmic derivative of the completed zeta function.
Its Hankel-square positivity for integer polynomials is equivalent to the Riemann
hypothesis.  Nevertheless, every compact-open neighbourhood of every admissible local
moment germ contains an exact normalized trace in the same closed finite-factor orbit
that fails a single finite Hankel-square test.  Hence no continuous local-germ reader
which is insensitive to finite Euler-factor changes can classify the RH-equivalent
positivity.

This does not decide the Riemann hypothesis.  It proves a sharp trilemma: an
RH-selecting source mechanism must retain genuinely global weighted-tail information,
be discontinuous in compact-open germ topology, or detect finite local-factor changes.
\end{abstract}

\section{Introduction}

Write
\[
 \xi(s)=\frac12s(s-1)\pi^{-s/2}\Gamma(s/2)\zeta(s),
 \qquad
 \Xi(z)=\xi\!\left(\frac12+iz\right).
 \tag{1.1}
\]
The Riemann hypothesis (RH) says that every zero of \(\Xi\) is real.  Many spectral,
trace-formula, moment, and Euler-product approaches seek a positivity object whose
nonnegativity forces this reality.  The hard step is not writing an equivalent
positivity criterion: it is deriving its sign from arithmetic source data without
already reading the zero divisor.

This paper isolates a previously hidden topological obstruction at the
functional-equation centre.  The obstruction is not specific to primes.  It is a
universal consequence of a translated Blaschke-divergence condition on a positive
frequency spectrum.

For a frequency \(\lambda>0\), put
\[
 W_\lambda(s)
  =\lambda\frac{e^{-\lambda s}}{(1-e^{-\lambda s})^2},
 \qquad
 U_\lambda(s)=W_\lambda(s)-W_\lambda(1-s).
 \tag{1.2}
\]
The reflection \(\FE s=1-s\) gives \(U_\lambda\circ\FE=-U_\lambda\).
The first theorem says that a broad class of frequency spectra makes the family
\(\{U_\lambda\}\) complete in every holomorphic \(\FE\)-odd critical-centre
fibre.  The exact hypothesis is stable under arbitrary finite deletion.

The linear statement is only the tangent law.  It upgrades to a nonlinear one.
For real \(b\) near zero define
\[
 \Phi_{\lambda,a}(s)
   =\lambda\frac{a e^{-\lambda s}}{1-ae^{-\lambda s}},
 \tag{1.3}
\]
and the reflected local increment
\[
 \Delta_{\lambda,b}(s)
 =\bigl[\Phi_{\lambda,1+b}(s)-\Phi_{\lambda,1}(s)\bigr]
  -\bigl[\Phi_{\lambda,1+b}(1-s)-\Phi_{\lambda,1}(1-s)\bigr].
 \tag{1.4}
\]
Whenever the perturbed denominators are nonzero on the carrier, this is the
logarithmic derivative of an invertible reflected local-factor modification.  Its
parameter derivative at zero is \(U_\lambda\).  Repetition of increasingly small
increments gives every real tangent multiple:
\[
 (N+1)\Delta_{\lambda,c/(N+1)}
   \longrightarrow cU_\lambda.
 \tag{1.5}
\]
Thus arbitrarily small genuine local factors generate a dense additive orbit.  Any
continuous Hausdorff-valued reader invariant under that group is constant on the
analytic fibre.

For the prime spectrum \(\lambda_p=\log p\), the required divergence is elementary:
for every fixed \(n\), eventually \((\log p)^n\leq p\), so Euler's theorem
\(\sum_p1/p=\infty\) implies the translated Blaschke divergence.  Finite prime
deletion changes nothing.  The prime instance is therefore a consequence of the
abstract spectrum theorem, not its proof mechanism.

The connection to RH is exact.  Near zero, the continued prime--archimedean source
trace equals \(\Xi'/\Xi\).  Its normalized divided difference has even Taylor
coefficients
\[
 M_n=\sum_{\rho}m_\rho
 \left[-\frac{1}{(\rho-\frac12)^2}\right]^{n+1},
 \tag{1.6}
\]
where the sum runs over the nontrivial zero divisor with multiplicity.  The quadratic
Hankel form
\[
 Q_\Xi(P)=\sum_\rho m_\rho\lambda_\rho P(\lambda_\rho)^2,
 \qquad
 \lambda_\rho=-\frac{1}{(\rho-\frac12)^2},
 \tag{1.7}
\]
is nonnegative for every integer polynomial \(P\) if and only if RH holds.  Hankel
reformulations of RH have substantial prior literature; no novelty is claimed for
the existence of such a criterion
\cite{Matiyasevich2007,Matiyasevich2017,BickelPascoeSargent2023}.

The new point is the collision of the RH-equivalent sign with the universal
finite-factor orbit.  Every compact-open neighbourhood of every even Schwarz-real
analytic moment germ contains another such germ which fails a single monomial-square
Hankel test, while its odd trace preimage lies in the closed pole-free local-factor
orbit.  Therefore local analytic continuation erases exactly the finite-factor data
that a continuous invariant sign classifier would need.

\subsection{Main results}

The paper proves four nested statements.

\begin{enumerate}[label=\textup{(\roman*)}]
\item A carrier-adapted completeness theorem needs only one sector order \(n\).
\item Requiring the translated Blaschke divergence for every \(n\) gives a
carrier-independent law on all compact positive-strip carriers and all
critical-centre discs.
\item The tangent law upgrades to a dense orbit generated only by pole-free
multiplicative increments; the maximal continuous invariant quotient on the
analytic fibre is trivial.
\item Prime logarithms satisfy the universal hypothesis cofinitely, and no
continuous finite-Euler-insensitive local classifier can decide the
RH-equivalent Hankel sign.
\end{enumerate}

The word ``universal'' refers to items (i)--(iii): the proof uses only the explicit
spectral divergence condition, not primes, an Euler product, a zero divisor, or RH.

\subsection{Relation to existing universality and completeness}

Classical exponential completeness asks when systems \(e^{\lambda z}\) are dense in
function spaces.  Its modern forms include the Beurling--Malliavin type problem,
M\"untz--Sz\'asz problems in holomorphic spaces, and geometric completeness criteria
\cite{Poltoratski2013,PelosoSalvatori2017,KhabibullinKudasheva2023}.
Our family is not a pure exponential system: it is the reflected
derivative-Lambert family in (1.2), and the proof must combine sectorial zero
uniqueness with Lambert--M\"obius recovery.

Voronin universality concerns vertical translates of \(\zeta\) in a substrip and
approximates nonvanishing holomorphic targets
\cite{Voronin1975,LamzouriLesterRadziwill2018}.  Generalized Euler-product
universality and finite Euler-product approximation are likewise distinct
\cite{Steuding2006,Gonek2007}.  Here the centre is fixed, the carrier crosses the
critical line, the reflection symmetry is built in, the target may vanish, and the
acting objects are genuine local-factor increments rather than vertical translates.

To the best of our knowledge, the cofinite reflected-Lambert completeness theorem,
its pole-free nonlinear orbit law, and the resulting RH local-classifier trilemma do
not appear in these literatures.  Priority is claimed only for this exact package.

\subsection{Scope}

The result does not rule out Hilbert--P\'olya, trace formulas, semilocal operator
limits, noncommutative geometry, or global Euler-product methods.  Such constructions
can evade the theorem by retaining a topology on the weighted prime tail which is
strictly finer than compact-open convergence of the continued germ, by reacting to
finite local changes, or by using a discontinuous/global invariant.  What is excluded
is the simultaneous use of:
\[
 \boxed{
 \text{compact-open local dependence}
 \;+\;
 \text{continuity}
 \;+\;
 \text{finite-factor invariance}
 }
 \tag{1.8}
\]
to classify the RH-equivalent trace-Hankel sign.

\section{Spectral hypothesis and main theorem}

\subsection{Translated sector-Blaschke completeness}

Let \(J\) be an arbitrary index set and
\(\Lambda=(\lambda_j)_{j\in J}\subset(0,\infty)\) an injective family.

\begin{definition}[Sector-Blaschke completeness]\label{def:SB}
Fix \(n\geq1\) and \(R_0\in\R\).  We say that \(\Lambda\) is
\((n,R_0)\)-sector-Blaschke complete if, for every \(R\geq R_0\),
\[
 \sum_{\lambda_j>R+1}
 \frac{1}{(\lambda_j-R)^n+1}=\infty.
 \tag{2.1}
\]
It is \emph{strongly sector-Blaschke complete above \(R_0\)} if
\textup{(2.1)} holds for every integer \(n\geq1\).
\end{definition}

All sums are over distinct frequencies.  The formulation permits uncountable \(J\),
although (2.1) then automatically contains a countable divergent subfamily.

\begin{proposition}[Cofinite stability]\label{prop:cofinite}
If \(\Lambda\) is strongly sector-Blaschke complete above \(R_0\), then
\(\Lambda\setminus F\) is strongly sector-Blaschke complete above \(R_0\) for every
finite \(F\subset J\).
\end{proposition}

\begin{proof}
Deleting finitely many nonnegative summands cannot turn a divergent sum into a
convergent one.  Injectivity is inherited.
\end{proof}

The strong condition is deliberately carrier-independent.  For a fixed carrier, one
sector order suffices; see Theorem~\ref{thm:carrier}.

\subsection{Critical-centre spaces and pole-free increments}

For \(0\leq r<1/2\), let
\[
 K_r=\left\{s\in\C:\left|s-\frac12\right|\leq r\right\},
 \qquad \FE(s)=1-s.
 \tag{2.2}
\]
Write \(C(K_r)\) for the complex Banach space with uniform norm.  Let
\(\cA^-_{\R}(K_r;r_1)\) be the restrictions to \(K_r\) of functions holomorphic on
\(|s-\frac12|<r_1\), continuous on its closure, and satisfying
\[
 f(1-s)=-f(s),\qquad f(\overline{s})=\overline{f(s)},
 \tag{2.3}
\]
where \(r<r_1<1/2\).

For \(\epsilon>0\), define \(\cG_{\Lambda}(r,\epsilon)\) as the additive group in
\(C(K_r)\) generated by all \(\Delta_{\lambda_j,b}\) in (1.4), with
\(b\in\R\), for which
\[
 |b|<\epsilon,
 \qquad
 \sup_{s\in K_r}|(1+b)e^{-\lambda_j s}|<1.
 \tag{2.4}
\]
Condition (2.4) makes both perturbed reflected local factors holomorphic and
invertible on \(K_r\).  Indeed reflection preserves \(K_r\), and the unperturbed
factor is safe because \(\Re s>0\).

\begin{theorem}[Universal reflected-Lambert orbit law]\label{thm:main}
Let \(\Lambda\subset(0,\infty)\) be injective and strongly sector-Blaschke complete
above some \(R_0>0\).  Let
\[
 0\leq r<r_1<\frac12,\qquad \epsilon>0.
 \tag{2.5}
\]
Then:
\begin{enumerate}[label=\textup{(\alph*)}]
\item the real span of \(\{U_{\lambda_j}:j\in J\}\) is dense in
\(\cA^-_{\R}(K_r;r_1)\) in the uniform norm;
\item every \(f\in\cA^-_{\R}(K_r;r_1)\) belongs to
\(\overline{\cG_{\Lambda}(r,\epsilon)}\);
\item if \(Y\) is Hausdorff and \(A:C(K_r)\to Y\) is continuous and satisfies
\[
 A(x+g)=A(x)\qquad(x\in C(K_r),\ g\in\cG_{\Lambda}(r,\epsilon)),
 \tag{2.6}
\]
then \(A(f)=A(0)\) for every \(f\in\cA^-_{\R}(K_r;r_1)\).
\end{enumerate}
All three conclusions remain true after deleting any finite set of frequencies.
\end{theorem}

Part (c) is the precise one-point quotient statement.  It concerns the analytic
\(\FE\)-odd Schwarz-real fibre, not all of \(C(K_r)\).

\begin{corollary}[Prime-logarithm law]\label{cor:prime}
For every finite set of primes \(S\), Theorem~\ref{thm:main} applies to
\[
 \Lambda_S=\{\log p:p\ \text{prime},\ p\notin S\}.
 \tag{2.7}
\]
Thus the cofinite prime reflected-Lambert tangents have dense real span, the
pole-free cofinite finite-Euler increment group has dense orbit, and every continuous
Hausdorff-valued invariant reader is constant on the analytic fibre.
\end{corollary}

\section{Sectorial zero uniqueness}

This section proves the analytic uniqueness mechanism behind the theorem.

\subsection{Cayley defects}

For \(n\geq1\) and \(R\in\R\), consider the translated sector
\[
 \Sigma_{n,R}
  =\left\{R+w:w\neq0,\ |\arg w|<\frac{\pi}{2n}\right\}.
 \tag{3.1}
\]
The map
\[
 \psi_{n,R}(z)
 =R+\left(\frac{1+z}{1-z}\right)^{1/n}
 \tag{3.2}
\]
maps the unit disc conformally to \(\Sigma_{n,R}\), with the principal root.
For \(\lambda>R+1\), its preimage is
\[
 a_{\lambda,n,R}
 =\frac{(\lambda-R)^n-1}{(\lambda-R)^n+1}\in(0,1).
 \tag{3.3}
\]
The exact radial defect is
\[
 1-|a_{\lambda,n,R}|
 =\frac{2}{(\lambda-R)^n+1}.
 \tag{3.4}
\]

\begin{lemma}[Translated Blaschke uniqueness]\label{lem:blaschke}
Suppose \(\Lambda\) is \((n,R_0)\)-sector-Blaschke complete.  Let \(H\) be a
function such that, for every \(R\geq R_0\), \(H\) is holomorphic and bounded on
\(\Sigma_{n,R}\).  If
\[
 H(\lambda_j)=0\qquad(\lambda_j>R_0+1),
 \tag{3.5}
\]
then
\[
 H(x)=0\qquad(x\geq R_0+1).
 \tag{3.6}
\]
\end{lemma}

\begin{proof}
Fix \(x\geq R_0+1\) and put \(R=x-1\).  If \(H(x)\neq0\), then
\[
 h(z)=H(\psi_{n,R}(z))
\]
is a bounded nonzero holomorphic function on \(\D\) whose distinct zeros include
every \(a_{\lambda_j,n,R}\) with \(\lambda_j>R+1\).  The zero set of a nonzero
bounded holomorphic disc function satisfies the Blaschke condition
\[
 \sum_j(1-|a_j|)<\infty
 \tag{3.7}
\]
\cite[Chapter II]{Garnett2007}.  Equations (3.4) and (2.1) contradict
\textup{(3.7)}.  Hence \(H(x)=0\).  Since \(x\) was arbitrary, (3.6) follows.
\end{proof}

The translated formulation is essential: it turns uniqueness at the spectral nodes
into vanishing at every point of a positive ray without invoking a density theorem
for the nodes.

\section{Lambert--M\"obius recovery on a compact carrier}

\subsection{Carrier-adapted theorem}

Let \(K\) be compact and \(s:K\to\C\) continuous.  Suppose
\[
 \Re s(k)\geq\delta>0,\qquad |\Im s(k)|\leq M.
 \tag{4.1}
\]
Define
\[
 W_\lambda(k)
 =\lambda\frac{e^{-\lambda s(k)}}{(1-e^{-\lambda s(k)})^2}.
 \tag{4.2}
\]

For a continuous linear functional \(L\in C(K)^*\), put
\[
 F(z)=L(e^{-zs}),\qquad
 H(z)=L\!\left(\frac{e^{-zs}}{(1-e^{-zs})^2}\right).
 \tag{4.3}
\]
The first is entire.  On every region where
\(\inf_{k\in K}\Re(zs(k))>0\), the geometric-derivative expansion gives
\[
 H(z)=\sum_{m\geq1}mF(mz),
 \tag{4.4}
\]
normally and locally uniformly.

\begin{theorem}[Carrier-adapted completeness]\label{thm:carrier}
Assume that \(\Lambda\) is \((n,R_0)\)-sector-Blaschke complete for one
\(n\geq1\), \(R_0\geq0\), and that for every \(R\geq R_0\) there is
\(d_R>0\) such that
\[
 \Re(zs(k))\geq d_R
 \quad(z\in\Sigma_{n,R},\ k\in K).
 \tag{4.5}
\]
If \(L(W_{\lambda_j})=0\) for every \(j\), then \(F\equiv0\).  Consequently
every coordinate monomial \(s^N\) belongs to the closed complex span of
\(\{W_{\lambda_j}\}\).
\end{theorem}

\begin{proof}
Since every \(\lambda_j>0\),
\[
 L(W_{\lambda_j})=\lambda_j H(\lambda_j)=0.
 \tag{4.6}
\]
Condition (4.5), compactness, and the uniform bound
\(|e^{-zs(k)}|<1\) imply that \(H\) is bounded and holomorphic on each
translated sector in Lemma~\ref{lem:blaschke}.  Hence
\[
 H(x)=0\qquad(x\geq R_0+1).
 \tag{4.7}
\]

Fix such an \(x\).  The positivity in (4.1) gives constants \(C>0\) and
\(0<q<1\) with
\[
 |F(mx)|\leq Cq^m.
 \tag{4.8}
\]
Thus all following sums are absolutely convergent.  M\"obius inversion of
\textup{(4.4)} yields
\[
\begin{aligned}
 \sum_{n\geq1}\mu(n)nH(nx)
 &=\sum_{n,m\geq1}\mu(n)nmF(mnx)\\
 &=\sum_{k\geq1}kF(kx)\sum_{n\mid k}\mu(n)
 =F(x).
\end{aligned}
\tag{4.9}
\]
Equation (4.7) therefore gives \(F(x)=0\) on a ray.  The identity theorem gives
\(F\equiv0\).

Differentiating at zero,
\[
 F^{(N)}(0)=(-1)^N L(s^N)=0
 \tag{4.10}
\]
for every \(N\).  If \(s^N\) were outside the closed span of the
\(W_{\lambda_j}\), Hahn--Banach would supply an \(L\) annihilating the family
but not \(s^N\), contradicting (4.10).
\end{proof}

\begin{remark}[Why the strong hypothesis is universal]
For a fixed positive rectangle (4.1), an integer \(n\) satisfying (4.5) exists:
choose the sector aperture small enough compared with \(\delta/M\).  Its value
depends on the carrier.  Strong sector-Blaschke completeness supplies every
integer \(n\), hence works simultaneously for every compact carrier in every
positive real strip.  This is the exact distinction between the carrier-adapted
and carrier-independent statements.
\end{remark}

\subsection{Reflection}

Let \(\tau:K\to K\) be continuous and suppose
\[
 s(\tau k)=1-s(k).
 \tag{4.11}
\]
Define \(P^-f=f-f\circ\tau\).  Then
\[
 P^-W_\lambda
 =W_\lambda-W_\lambda\circ\tau.
 \tag{4.12}
\]

\begin{proposition}[Reflected polynomial completeness]\label{prop:poly}
Under the strong sector-Blaschke hypothesis and (4.1), every reflected polynomial
\[
 P(s)-P(1-s)
 \tag{4.13}
\]
belongs to the closed complex span of
\(\{W_{\lambda_j}(s)-W_{\lambda_j}(1-s)\}\).
\end{proposition}

\begin{proof}
If a continuous functional annihilates all reflected tangents, compose it with
\(P^-\).  The resulting functional annihilates every unreflected \(W_{\lambda_j}\).
Theorem~\ref{thm:carrier} makes it vanish on every coordinate monomial, hence on
every polynomial.  Hahn--Banach gives (4.13).
\end{proof}

\section{Holomorphic and Schwarz-real completeness}

\subsection{Passage from polynomials to analytic germs}

Fix \(0\leq r<r_1<1/2\).  Let \(f\) be holomorphic on
\(|s-\frac12|<r_1\), continuous on the closure, and \(\FE\)-odd on \(K_r\).
Its Taylor polynomials \(P_N\) about \(1/2\) converge uniformly on \(K_r\).
Oddification preserves convergence:
\[
 \frac12\bigl(P_N(s)-P_N(1-s)\bigr)\longrightarrow f(s)
 \quad\text{uniformly on }K_r.
 \tag{5.1}
\]
Proposition~\ref{prop:poly} therefore proves complex density of the reflected
Lambert family.

To descend to the real form, define Schwarz conjugation
\[
 f^\#(s)=\overline{f(\overline{s})}.
 \tag{5.2}
\]
Every \(U_{\lambda_j}\) is fixed by \(\#\), because \(\lambda_j\) is real.  If \(f=f^\#\)
and complex linear combinations of \(U_{\lambda_j}\) approximate \(f\), applying the
continuous projection
\[
 \operatorname{Re}_\# h=\frac{h+h^\#}{2}
 \tag{5.3}
\]
replaces them by real linear combinations with the same limit.  This proves
Theorem~\ref{thm:main}(a).

\subsection{A Fr\'echet-space formulation}

Let
\[
 D_{1/2}=\left\{s:\left|s-\frac12\right|<\frac12\right\}.
 \tag{5.4}
\]
Endow the Schwarz-real \(\FE\)-odd holomorphic space
\(\cH^-_{\R}(D_{1/2})\) with compact-open topology.

\begin{corollary}[Germ-space completeness]\label{cor:frechet}
For every strongly sector-Blaschke complete positive spectrum, the real span of
\(\{U_{\lambda_j}\}\) is dense in \(\cH^-_{\R}(D_{1/2})\).  The same holds after
arbitrary finite spectral deletion.
\end{corollary}

\begin{proof}
A basic compact-open neighbourhood is controlled on a compact subset of
\(D_{1/2}\), hence on some \(K_r\).  Choose \(r<r_1<1/2\) and apply
Theorem~\ref{thm:main}(a).
\end{proof}

\section{The pole-free nonlinear orbit}

\subsection{Tangent identification}

Let \(q_\lambda(s)=e^{-\lambda s}\).  On \(K_r\),
\(\|q_\lambda\|_\infty<1\).  The map
\[
 a\longmapsto \Phi_{\lambda,a}
  =\lambda\frac{a q_\lambda}{1-aq_\lambda}
 \tag{6.1}
\]
is holomorphic in a neighbourhood of \(a=1\) as a map into the Banach algebra
\(C(K_r)\), and
\[
 \left.\frac{d}{da}\Phi_{\lambda,a}\right|_{a=1}
 =\lambda\frac{q_\lambda}{(1-q_\lambda)^2}
 =W_\lambda.
 \tag{6.2}
\]
After reflection,
\[
 \left.\frac{d}{db}\Delta_{\lambda,b}\right|_{b=0}=U_\lambda.
 \tag{6.3}
\]
Real \(b\) preserves Schwarz reality.

Pole-freeness is not an auxiliary convention.  Put \(q=\|q_\lambda\|_\infty<1\).
If
\[
 |b|<\eta_\lambda:=\frac{1-q}{1+q},
 \tag{6.4}
\]
then
\[
 \|(1+b)q_\lambda\|_\infty
 \leq(1+|b|)q<1.
 \tag{6.5}
\]
Thus each individual tangent has a genuine pole-free parameter neighbourhood,
even if the admissible radius is not uniform in \(\lambda\).

\subsection{Repetition lemma}

\begin{lemma}[Local increments generate the tangent]\label{lem:repeat}
Let \(E\) be a real Banach space and \(\Delta:(-\eta,\eta)\to E\) be
differentiable at zero, with \(\Delta(0)=0\) and \(\Delta'(0)=u\).  For every
\(c\in\R\),
\[
 (N+1)\Delta\!\left(\frac{c}{N+1}\right)\longrightarrow cu.
 \tag{6.6}
\]
Consequently \(cu\) lies in the closure of the additive group generated by
\(\Delta((-\eta,\eta))\).
\end{lemma}

\begin{proof}
The definition of the derivative gives
\[
 \frac{\Delta(t)}{t}\longrightarrow u\qquad(t\to0).
 \]
Set \(t=c/(N+1)\).  For \(c=0\) the claim is immediate; otherwise
\[
 (N+1)\Delta\!\left(\frac{c}{N+1}\right)
 =c\,\frac{\Delta(t)}{t}\longrightarrow cu.
\]
Each term on the left is a natural-number repetition of one allowed group
generator.
\end{proof}

For every prescribed \(\epsilon>0\), the parameters \(c/(N+1)\) eventually satisfy
both \(|b|<\epsilon\) and the pole-free bound (6.4).  Applying
Lemma~\ref{lem:repeat} frequency by frequency shows that the closure of
\(\cG_\Lambda(r,\epsilon)\) contains the real span of all \(U_{\lambda_j}\).
Theorem~\ref{thm:main}(b) follows from linear density.

\subsection{One-point continuous quotient}

Let \(A:C(K_r)\to Y\) satisfy (2.6), with \(Y\) Hausdorff.  The fibre
\[
 E_0=\{x:A(x)=A(0)\}
\tag{6.7}
\]
is closed by continuity and Hausdorff separation.  It contains
\(\cG_\Lambda(r,\epsilon)\), hence its closure.  Part (b) gives
\(\cA^-_\R(K_r;r_1)\subset E_0\), proving Theorem~\ref{thm:main}(c).

\begin{remark}[Multiplicative meaning]
Put
\[
 V_{\lambda,b}(s)
 =\frac{1-(1+b)e^{-\lambda s}}{1-e^{-\lambda s}}.
 \tag{6.8}
\]
On the pole-free parameter locus,
\[
 \Delta_{\lambda,b}(s)
 =\frac{d}{ds}\log\!\bigl(V_{\lambda,b}(s)V_{\lambda,b}(1-s)\bigr).
 \tag{6.9}
\]
Finite sums and negatives in \(\cG_\Lambda\) are therefore logarithmic
derivatives of finite products and ratios of invertible reflected local factors.
Repetition in (6.6) produces a power of a local factor; it is not asserted to remain
in the raw degree-one parameter locus.
\end{remark}

\section{Prime logarithms}

\begin{proposition}[Cofinite strong completeness of prime logs]\label{prop:primelog}
For every finite prime set \(S\), every \(R_0\geq0\), every \(n\geq1\), and
every \(R\geq R_0\),
\[
 \sum_{\substack{p\notin S\\ \log p>R+1}}
 \frac{1}{(\log p-R)^n+1}=\infty.
 \tag{7.1}
\]
\end{proposition}

\begin{proof}
Fix \(n\) and \(R\).  For all sufficiently large primes,
\[
 (\log p-R)^n+1\leq 2(\log p)^n\leq 2p.
 \tag{7.2}
\]
Hence the summand in (7.1) is at least \(1/(2p)\) eventually.  Euler's theorem
\[
 \sum_p\frac1p=\infty
 \tag{7.3}
\]
implies (7.1), and deleting finitely many terms does not affect divergence.
The map \(p\mapsto\log p\) is injective and positive.
\end{proof}

Corollary~\ref{cor:prime} is immediate.  No prime number theorem, zero-free region,
or assertion about zeta zeros enters the proof.

\section{Normalized trace germs and Hankel positivity}

\subsection{Odd-to-even normalization}

Let \(T\) be an odd analytic germ at zero with \(T(0)=0\).  Define
\[
 \mathcal N T(z)=-\frac{T(z)}{z},
\tag{8.1}
\]
where the value at zero is the removable value \(-T'(0)\).  Then
\(\mathcal NT\) is even.  Conversely, for an even germ \(g\),
\[
 \mathcal I g(z)=-zg(z)
\tag{8.2}
\]
is odd and
\[
 \mathcal N\mathcal I g=g.
\tag{8.3}
\]

In the critical coordinate \(s\), define
\[
 (\cT g)(s)
 =-\left(s-\frac12\right)
    g\!\left(-i\left(s-\frac12\right)\right).
\tag{8.4}
\]
If \(g\) is even and Schwarz-real, then \(\cT g\) is \(\FE\)-odd and
Schwarz-real.  Rotating back,
\[
 -i(\cT g)\!\left(\frac12+iz\right)=-zg(z),
\tag{8.5}
\]
so normalized division recovers \(g\) exactly.

\begin{corollary}[Moment-germ orbit transport]\label{cor:momentorbit}
Under the hypotheses of Theorem~\ref{thm:main}, every even Schwarz-real analytic
germ \(g\) on \(|z|<r_1\) is the exact normalized germ of the trace \(\cT g\), and
\[
 \cT g|_{K_r}\in\overline{\cG_\Lambda(r,\epsilon)}.
\tag{8.6}
\]
This remains true after any finite spectral deletion.
\end{corollary}

\begin{proof}
Equation (8.3) gives exact normalization.  The symmetries of (8.4) put
\(\cT g\) in the fibre of Theorem~\ref{thm:main}(b).
\end{proof}

\subsection{The analytic Hankel cone has no local thickness}

For an even Schwarz-real analytic germ
\[
 g(z)=\sum_{n\geq0}m_n(g)z^{2n},
\qquad m_n(g)\in\R,
\tag{8.7}
\]
and a real polynomial \(P(X)=\sum_jc_jX^j\), define
\[
 Q_g(P)=\sum_{i,j}c_ic_jm_{i+j}(g).
\tag{8.8}
\]
Let
\[
 \cP_H=\{g:Q_g(P)\geq0\ \text{for every }P\in\R[X]\}.
\tag{8.9}
\]

\begin{proposition}[Remote-coefficient instability]\label{prop:instability}
Let \(0\leq r<r_1<1/2\).  For every even Schwarz-real \(g\), every
\(\eta>0\), and every uniform neighbourhood on \(K_r\), there is an even
Schwarz-real analytic \(\widetilde g\) such that
\[
 \sup_{|z|\leq r}|\widetilde g(z)-g(z)|<\eta,
 \qquad
 \widetilde g\notin\cP_H.
\tag{8.10}
\]
The failure is witnessed by one monomial square.
\end{proposition}

\begin{proof}
Choose \(B\geq0\) with \(|g(z)|\leq B\) on \(|z|=r_1\).  Cauchy's estimate gives
\[
 |m_{2N}(g)|\leq Br_1^{-4N}.
\tag{8.11}
\]
Set
\[
 g_N(z)=g(z)-\bigl(m_{2N}(g)+1\bigr)z^{4N}.
\tag{8.12}
\]
Then \(g_N\) preserves both symmetries and its \(z^{4N}\)-coefficient is exactly
\(-1\).  Hence
\[
 Q_{g_N}(X^N)=m_{2N}(g_N)=-1.
\tag{8.13}
\]
On \(|z|\leq r\),
\[
\begin{aligned}
 |g_N(z)-g(z)|
 &\leq\bigl(|m_{2N}(g)|+1\bigr)r^{4N}\\
 &\leq B(r/r_1)^{4N}+r^{4N}\longrightarrow0,
\end{aligned}
\tag{8.14}
\]
because \(r<r_1\) and \(r<1\).  Choose \(N\) large enough for (8.10).
\end{proof}

Combining Proposition~\ref{prop:instability} with
Corollary~\ref{cor:momentorbit} gives more than abstract empty interior.

\begin{corollary}[Orbit-realized sign failure]\label{cor:orbitfailure}
Every compact-disc neighbourhood of every even Schwarz-real analytic moment germ
contains an exact even germ \(\widetilde g\notin\cP_H\) whose critical trace
\(\cT\widetilde g\) belongs to
\(\overline{\cG_\Lambda(r,\epsilon)}\).
\end{corollary}

\section{The completed zeta source}

\subsection{Prime--archimedean trace identity}

Put
\[
 A(s)=\frac12s(s-1)\pi^{-s/2}\Gamma(s/2),
 \qquad \xi(s)=A(s)\zeta(s).
\tag{9.1}
\]
On \(\Re s>1\),
\[
 -\frac{\zeta'}{\zeta}(s)
 =\sum_{n\geq1}\frac{\Lambda(n)}{n^s}.
\tag{9.2}
\]
Its meromorphic continuation and the fixed archimedean correction define near
\(z=0\) the source trace
\[
\begin{aligned}
 T_\zeta(z)
 &=-i\left(
   -\frac{\zeta'}{\zeta}\!\left(\frac12+iz\right)
   -\frac{A'}{A}\!\left(\frac12+iz\right)\right)\\
 &=\frac{\Xi'(z)}{\Xi(z)}.
\end{aligned}
\tag{9.3}
\]
The second equality is the logarithmic derivative of (9.1) under
\(s=\frac12+iz\).  Since \(\Xi\) is even and \(\Xi(0)\neq0\), \(T_\zeta\) is an
odd analytic germ and \(T_\zeta(0)=0\).

Define the normalized prime--archimedean moment germ
\[
 G_\zeta(z)=-\frac{T_\zeta(z)}{z},
\tag{9.4}
\]
with its removable value at zero.

\subsection{Reciprocal-square divisor moments}

For each nontrivial zero \(\rho\) of \(\xi\), counted with multiplicity \(m_\rho\),
put
\[
 \lambda_\rho=-\frac{1}{(\rho-\frac12)^2}.
\tag{9.5}
\]
The centered inverse-square zero sum is absolutely convergent.  The logarithmic
derivative of the Hadamard product, paired by the functional equation, gives for
sufficiently small \(z\) \cite[Chapter~2]{Titchmarsh1986}
\[
 G_\zeta(z)
 =\sum_\rho m_\rho\frac{\lambda_\rho}{1-z^2\lambda_\rho}
 =\sum_{n\geq0}
   \left(\sum_\rho m_\rho\lambda_\rho^{n+1}\right)z^{2n}.
\tag{9.6}
\]
Thus
\[
 m_n(G_\zeta)=M_n
 :=\sum_\rho m_\rho\lambda_\rho^{n+1}.
\tag{9.7}
\]
Conjugation symmetry of the divisor makes every \(M_n\) real unconditionally.

For \(P\in\R[X]\), absolute convergence permits
\[
\begin{aligned}
 Q_{G_\zeta}(P)
 &=\sum_{i,j}c_ic_jM_{i+j}\\
 &=\sum_\rho m_\rho\lambda_\rho P(\lambda_\rho)^2.
\end{aligned}
\tag{9.8}
\]

\begin{lemma}[Trace-class polynomial detector]\label{lem:detector}
Let \((\lambda_j)_{j\in J}\subset\C\setminus\{0\}\) have positive integer
weights \(w_j\), and suppose
\[
 \sum_jw_j|\lambda_j|<\infty.
\tag{9.9}
\]
If
\[
 \Re\sum_jw_j\lambda_jP(\lambda_j)^2\geq0
 \qquad(P\in\R[X]),
\tag{9.10}
\]
then every \(\lambda_j\) is real.
\end{lemma}

\begin{proof}
Assume that \(\alpha=\lambda_{j_0}\) is nonreal.  Choose \(c\in\C\) such that
\[
 \Re(\alpha c^2)<0.
\tag{9.11}
\]
Such a phase always exists: \(c=1\) works if \(\Re\alpha<0\), \(c=i\) if
\(\Re\alpha>0\), and \(c=1\pm i\), with the appropriate sign, if
\(\Re\alpha=0\).

By (9.9), only finitely many indices satisfy
\(|\lambda_j|\geq|\alpha|\).  Let \(F\) be the finite set of their spectral
values.  There is a real polynomial \(S\) which vanishes on
\(F\setminus\{\alpha,\overline\alpha\}\) and does not vanish at \(\alpha\):
one may take a product of the real quadratics
\[
 (X-\beta)(X-\overline\beta)
 =X^2-2\Re(\beta)X+|\beta|^2
\tag{9.12}
\]
over the unwanted conjugation orbits.

Put \(u=\alpha/|\alpha|\).  For every \(N\), choose a real affine polynomial
\(L_N\) satisfying
\[
 L_N(\alpha)=\frac{c}{S(\alpha)u^N}.
\tag{9.13}
\]
This is possible because \(\Im\alpha\neq0\): for any \(d\in\C\),
\[
 L_d(X)=
 \Re d-\frac{\Im d}{\Im\alpha}\Re\alpha
 +\frac{\Im d}{\Im\alpha}X
\tag{9.14}
\]
has real coefficients and \(L_d(\alpha)=d\).  The coefficients of \(L_N\)
are bounded independently of \(N\), since \(|u|=1\).  Define
\[
 P_N(X)=S(X)L_N(X)\left(\frac{X}{|\alpha|}\right)^N.
\tag{9.15}
\]
Then
\[
\begin{array}{lll}
P_N(\alpha)=c,&
P_N(\overline\alpha)=\overline c,&
P_N(\beta)=0\quad
  (\beta\in F\setminus\{\alpha,\overline\alpha\}).
\end{array}
\tag{9.16}
\]
Moreover, \(P_N\) is bounded uniformly in \(N\) on
\(|z|\leq|\alpha|\), and \(P_N(z)\to0\) whenever \(|z|<|\alpha|\).

The summable majorant
\[
 C^2w_j|\lambda_j|
\tag{9.17}
\]
therefore permits dominated convergence in (9.10).  All outer spectral
values except \(\alpha,\overline\alpha\) vanish identically; every inner value
vanishes in the limit.  The limiting real part is a positive sum of copies of
\(\Re(\alpha c^2)\), one for each occurrence of \(\alpha\) or
\(\overline\alpha\).  It is strictly negative by (9.11), contradicting
(9.10).  Hence no nonreal \(\lambda_j\) exists.
\end{proof}

\begin{theorem}[Countable source-Hankel criterion]\label{thm:hankelRH}
The following are equivalent:
\begin{enumerate}[label=\textup{(\roman*)}]
\item the Riemann hypothesis;
\item \(Q_{G_\zeta}(P)\geq0\) for every \(P\in\R[X]\);
\item \(Q_{G_\zeta}(P)\geq0\) for every \(P\in\mathbb Z[X]\).
\end{enumerate}
\end{theorem}

\begin{proof}
If RH holds, write \(\rho=\frac12+i\gamma\) with \(\gamma\in\R\setminus\{0\}\).
Then \(\lambda_\rho=\gamma^{-2}>0\), and every summand in (9.8) is
nonnegative.  Thus (i) implies (ii), and (ii) implies (iii).

For the converse from (ii), the reciprocal-square spectrum is trace class and has
positive integer weights, so Lemma~\ref{lem:detector} forces every
\(\lambda_\rho\) to be real.
Direct calculation gives
\[
 \Im\lambda_\rho
 =\frac{2(\Re\rho-\frac12)\Im\rho}
        {|\rho-\frac12|^4}.
\tag{9.18}
\]
Every nontrivial zeta zero has nonzero imaginary part.  Thus
\(\lambda_\rho\in\R\) is equivalent to \(\Re\rho=1/2\), and RH follows.

It remains to show that integer tests suffice.  At each fixed degree, \(Q\) is a
continuous quadratic form in finitely many real coefficients.  Rational coefficient
vectors are dense.  Clearing a common positive denominator turns a rational
polynomial into an integer polynomial, while quadratic homogeneity preserves sign.
Thus (iii) implies (ii).
\end{proof}

\begin{remark}
The separation step in the preceding proof can be made explicit by a phase-corrected
real interpolation polynomial on a finite conjugation-stable spectral set, raised to
a high power so that the trace-class tail is negligible.  The argument uses no zero
enumeration and no simplicity assumption.
\end{remark}

\section{The local-source barrier}

We can now state the RH consequence without conflating an equivalent criterion with
a proof of its sign.

\begin{theorem}[No continuous finite-factor-invariant local classifier]
\label{thm:noclassifier}
Let \(\Lambda\) satisfy Theorem~\ref{thm:main}, and let \(0\leq r<r_1<1/2\).
There is no continuous map \(A:C(K_r)\to Y\) to a Hausdorff space which is
invariant under \(\cG_\Lambda(r,\epsilon)\) and satisfies, for every analytic
Schwarz-real even germ \(g\),
\[
 A(\cT g|_{K_r})=A(0)
 \quad\Longleftrightarrow\quad
 Q_g(P)\geq0\ \text{for every }P\in\R[X].
\]
\end{theorem}

\begin{proof}
Every trace \(\cT g\) lies in the closed orbit of zero by
Corollary~\ref{cor:momentorbit}.  Every continuous invariant map therefore takes the
same value on \(\cT g\) and on zero by Theorem~\ref{thm:main}(c).  But the property
is nonconstant: the zero germ satisfies every Hankel inequality, whereas
\(g(z)=-1\) fails the test \(P=1\).  No such map can classify it.
\end{proof}

The explicit remote perturbation strengthens the contradiction.

\begin{theorem}[Local RH-source trilemma]\label{thm:trilemma}
For the prime spectrum \(\lambda_p=\log p\), after deletion of any finite prime set:
\begin{enumerate}[label=\textup{(\alph*)}]
\item the pole-free finite-Euler increment orbit is dense in every analytic
Schwarz-real \(\FE\)-odd critical-centre fibre;
\item normalized prime--archimedean source-Hankel positivity is equivalent to RH;
\item every compact-open neighbourhood in the normalized moment-germ fibre contains
an exact trace from the same closed finite-Euler orbit which fails a finite
Hankel-square test.
\end{enumerate}
Consequently, any source mechanism which establishes the RH-equivalent sign must
violate at least one of:
\[
\begin{array}{ll}
\textup{(L)}&\text{dependence only on a compact-open critical-centre germ},\\
\textup{(C)}&\text{continuity in that germ topology},\\
\textup{(F)}&\text{invariance under finite pole-free Euler-factor changes}.
\end{array}
\tag{10.1}
\]
\end{theorem}

\begin{proof}
Part (a) is Corollary~\ref{cor:prime}; part (b) is
Theorem~\ref{thm:hankelRH}; part (c) is
Corollary~\ref{cor:orbitfailure} specialized to prime logs.  If a mechanism
satisfied (L), (C), and (F), it would induce the classifier excluded by
Theorem~\ref{thm:noclassifier}.
\end{proof}

\subsection{What the trilemma decides}

Theorem~\ref{thm:trilemma} answers the local-source question completely.  The
continued analytic germ, even with all Taylor coefficients and exact
functional-equation and Schwarz symmetries, cannot carry a nonconstant continuous
finite-factor-invariant sign.  This is not a finite-jet limitation: density holds in
compact-open holomorphic topology, and the failure can be placed in every
neighbourhood.

The surviving source information is therefore global and arithmetic.  In the
half-plane of absolute convergence, (9.2) remembers the full weighted coefficient
graph
\[
 \bigl\{(\log p^k,\log p):p\ \text{prime},\ k\geq1\bigr\}.
\tag{10.2}
\]
Analytic continuation to a local germ forgets this presentation.  A viable
source-side RH proof must retain an incidence relation between that weighted
all-prime tail and the archimedean completion before passing to the germ quotient,
or else use one of the two other escape routes in (10.1).  The trilemma identifies
the required category; it does not assume the missing positivity.

\section{Sharpness and limitations}

\subsection{Sharpness of the analytic input}

The divergence in Definition~\ref{def:SB} is exactly the negation of the Blaschke
summability forced on zeros of a bounded analytic function after the sector-to-disc
map.  Thus it is sharp for the zero-uniqueness step used here.  It is a sufficient,
not claimed necessary, condition for reflected-Lambert completeness: a different
transform argument could conceivably prove completeness for sparser spectra.

The universal theorem uses every integer sector order only because the carrier is
arbitrary.  The carrier-adapted Theorem~\ref{thm:carrier} uses one order and is the
minimal statement delivered by this method.

\subsection{Topology}

Compact-open convergence is natural for holomorphic germs and controls all
derivatives on strictly smaller discs by Cauchy's formula.  It does not control the
arithmetic presentation of a continued function by a weighted Euler tail.  The
theorem is therefore not a paradox: distinct global sources can be arbitrarily close
or even indistinguishable to continuous local readers while remaining different in
a stronger source topology.

\subsection{No global spectral no-go}

The theorem does not prove that a famous route can never prove RH in every form.
In particular it does not exclude:
\begin{itemize}
\item a self-adjoint operator whose domain or boundary conditions depend on the full
global source;
\item a semilocal trace formula with a topology retaining place-by-place weights;
\item a discontinuous index, winding number, divisor count, or boundary invariant;
\item a construction deliberately sensitive to one or more finite Euler factors.
\end{itemize}
It excludes precisely the conjunction (L)+(C)+(F), and nothing broader.

\subsection{Comparison with zeta universality}

Voronin's theorem makes vertical translates of one global function approximate
nonvanishing targets on compacta contained in \(1/2<\Re s<1\).  Our theorem fixes the
centre, crosses the critical line, enforces \(s\leftrightarrow1-s\), permits zero
targets, and uses an additive group of actual local-factor logarithmic derivatives.
Neither theorem implies the other.

\subsection{Comparison with finite Euler products}

Finite Euler products are often used to approximate \(\zeta\) in regions depending
on height and, in strong forms, on RH \cite{Gonek2007}.  The present result is
unconditional and makes no approximation claim for \(\zeta\) itself.  It classifies
the orbit closure of local logarithmic source changes inside a fixed
critical-centre analytic fibre.

\section{Conclusion}

The reflected derivative-Lambert family obeys a universal cofinite completeness law.
Its hypothesis is spectral, translated, and carrier-independent; prime logarithms
satisfy it by Euler's reciprocal-prime divergence.  Linear density upgrades to a
dense nonlinear orbit generated entirely by arbitrarily small pole-free local
multiplicative modifications.  On the analytic functional-equation-odd
Schwarz-real fibre, the maximal continuous invariant quotient is therefore one
point.

At the Riemann source, normalized trace-Hankel positivity is exactly equivalent to
RH, yet that sign is nowhere protected by the compact-open local orbit: every
neighbourhood contains an exact orbit-realized failure.  The result closes the
continuous local-germ, finite-factor-insensitive route.  What remains viable is
sharply determined: a successful source proof must retain global weighted-tail
incidence, respond to finite local changes, or use a discontinuous invariant.

\begin{thebibliography}{99}

\bibitem{BickelPascoeSargent2023}
K.~Bickel, J.~E.~Pascoe, and M.~Sargent,
\newblock Zero-free regions near a line,
\newblock \emph{Math. Z.} \textbf{303} (2023), article 14.
\newblock \href{https://doi.org/10.1007/s00209-022-03149-5}
{doi:10.1007/s00209-022-03149-5}.

\bibitem{Garnett2007}
J.~B.~Garnett,
\newblock \emph{Bounded Analytic Functions},
\newblock revised first edition, Springer, New York, 2007.

\bibitem{Gonek2007}
S.~M.~Gonek,
\newblock Finite Euler products and the Riemann hypothesis,
\newblock \emph{Trans. Amer. Math. Soc.} \textbf{364} (2012), 2157--2191.
\newblock \href{https://arxiv.org/abs/0704.3448}{arXiv:0704.3448}.

\bibitem{KhabibullinKudasheva2023}
B.~N.~Khabibullin and E.~G.~Kudasheva,
\newblock Completeness of exponential systems in spaces of functions on a
compact or domain and their area,
\newblock \href{https://arxiv.org/abs/2306.15899}{arXiv:2306.15899}, 2023.

\bibitem{LamzouriLesterRadziwill2018}
Y.~Lamzouri, S.~Lester, and M.~Radziwi{\l}{\l},
\newblock An effective universality theorem for the Riemann zeta-function,
\newblock \emph{Comment. Math. Helv.} \textbf{93} (2018), 709--765.
\newblock \href{https://arxiv.org/abs/1611.10325}{arXiv:1611.10325}.

\bibitem{Matiyasevich2007}
Y.~Matiyasevich,
\newblock Hidden life of Riemann's zeta function 2. Electrons and trains,
\newblock \href{https://arxiv.org/abs/0709.0028}{arXiv:0709.0028}, 2007.

\bibitem{Matiyasevich2017}
Y.~V.~Matiyasevich,
\newblock Riemann's hypothesis in terms of the eigenvalues of special Hankel
matrices,
\newblock \emph{Proc. Steklov Inst. Math.} \textbf{296} Suppl.~2 (2017),
78--91.
\newblock \href{https://doi.org/10.1134/S0081543817030099}
{doi:10.1134/S0081543817030099}.

\bibitem{PelosoSalvatori2017}
M.~M.~Peloso and M.~Salvatori,
\newblock Functions of exponential growth in a half-plane, sets of uniqueness
and the M\"untz--Sz\'asz problem for the Bergman space,
\newblock \emph{J. Geom. Anal.} \textbf{27} (2017), 2570--2599.
\newblock \href{https://doi.org/10.1007/s12220-017-9773-8}
{doi:10.1007/s12220-017-9773-8}.

\bibitem{Poltoratski2013}
A.~Poltoratski,
\newblock A problem on completeness of exponentials,
\newblock \emph{Ann. of Math.} (2) \textbf{178} (2013), 983--1016.
\newblock \href{https://doi.org/10.4007/annals.2013.178.3.4}
{doi:10.4007/annals.2013.178.3.4}.

\bibitem{Steuding2006}
R.~Steuding,
\newblock Universality for generalized Euler products,
\newblock \emph{Analysis (Munich)} \textbf{26} (2006), 337--345.
\newblock \href{https://doi.org/10.1524/anly.2006.26.99.337}
{doi:10.1524/anly.2006.26.99.337}.

\bibitem{Titchmarsh1986}
E.~C.~Titchmarsh,
\newblock \emph{The Theory of the Riemann Zeta-Function},
\newblock second edition, revised by D.~R.~Heath-Brown, Oxford University
Press, Oxford, 1986.

\bibitem{Voronin1975}
S.~M.~Voronin,
\newblock Theorem on the universality of the Riemann zeta-function,
\newblock \emph{Math. USSR-Izv.} \textbf{9} (1975), 443--453.
\newblock \href{https://doi.org/10.1070/IM1975v009n03ABEH001485}
{doi:10.1070/IM1975v009n03ABEH001485}.

\end{thebibliography}

\end{document}
